\section{Introduction} \label{section:introduction}
Reinforcement learning (RL) has experienced significant advancements in the past years, enabling robots to achieve human like performance on complex tasks. A popular category of reinforcement learning algorithms include actor-critic reinforcement learning, which has shown pivotal in training efficiently in challenging environments. These algorithms have the potential of being more sample efficient, thanks to their parallel training. Concurrently, neuromorphic algorithms, such as spiking neural networks (SNN) have garnered attention in the field of deep learning. These bio-inspired algorithms have demonstrated utility in handling temporal data and offer high energy efficiency when run on specialized hardware. Combining the possibility to learn from experience from reinforcement learning with the energy efficient characteristics of spiking neural networks holds great opportunity in the field of small mobile robotics, where computational resources are constrained.

Current research in the field of spiking based reinforcement learning algorithms is largely based on the principle of deep Q learning (DQN). While showing impressive results, this method relies on the use of a technique called experience replay to reduce the bias during training. This experience therefore needs to be gathered before the training can start. Actor-critic reinforcement learning on the other hand, makes use of multiple workers operating in parallel, eliminating the need for experience replay. This opens doors to learning on resource constrained systems, where the overhead required for storing past experiences and training in batches of experience can be detrimental for the performance of the system.

Considering the advantages of the actor-critic algorithms and spiking neural networks, this article explores the challenges and opportunities of combining the two. Firstly, the related work in the field of spiking based reinforcement learning is discussed in \autoref{section:lit_review}. Going over the general use cases where these bio-inspired algorithms shine and the implementation details of each work. Next, the methods used for this article are discussed in \autoref{section:methods}. Then, \autoref{section:results} presents the results. In \autoref{section:recs}, the results are discussed and directions for future work are given. Lastly, the article is concluded in \autoref{section:Conc_Recc}.

